% Some basic packages
\usepackage[utf8]{inputenc}
\usepackage[T1]{fontenc}
\usepackage{textcomp}
\usepackage[english]{babel}
\usepackage{url}
\usepackage{graphicx}
\usepackage{float}
\usepackage{booktabs}
\usepackage{enumitem}
\usepackage{listings}
\usepackage{dcolumn}
\newcolumntype{2}{D{.}{}{2.0}}
\usepackage{xcolor}

\definecolor{codegreen}{rgb}{0,0.6,0}
\definecolor{codegray}{rgb}{0.5,0.5,0.5}
\definecolor{codepurple}{rgb}{0.58,0,0.82}
\definecolor{backcolour}{rgb}{0.95,0.95,0.92}

\lstdefinestyle{mystyle}{
    backgroundcolor=\color{backcolour},   
    commentstyle=\color{codegreen},
    keywordstyle=\color{magenta},
    numberstyle=\tiny\color{codegray},
    stringstyle=\color{codepurple},
    basicstyle=\ttfamily\footnotesize,
    breakatwhitespace=false,         
    breaklines=true,                 
    captionpos=b,                    
    keepspaces=true,                 
    numbers=left,                    
    numbersep=5pt,                  
    showspaces=false,                
    showstringspaces=false,
    showtabs=false,                  
    tabsize=2
}

\lstset{style=mystyle}

\pdfminorversion=7

\usepackage[
    backend=biber,
    style=authoryear-icomp,
    sortlocale=auto,
    natbib=true,
    url=false, 
    doi=true,
    eprint=false
]{biblatex}
% \addbibresource{example.bib}

\usepackage{hyperref}
\hypersetup{
	colorlinks=true,
}

% Don't indent paragraphs, leave some space between them
\usepackage{parskip}

% Hide page number when page is empty
\usepackage{emptypage}
\usepackage{subcaption}
\usepackage{multicol}

% Other font I sometimes use.
% \usepackage{cmbright}

% Math stuff
\usepackage{amsmath, amsfonts, mathtools, amsthm, amssymb}
% Fancy script capitals
\usepackage{mathrsfs}
\usepackage{cancel}
% Bold math
\usepackage{bm}
% Some shortcuts
\newcommand\N{\ensuremath{\mathbb{N}}}
\newcommand\R{\ensuremath{\mathbb{R}}}
\newcommand\Z{\ensuremath{\mathbb{Z}}}
%\renewcommand\O{\ensuremath{\emptyset}}
%\newcommand\Q{\ensuremath{\mathbb{Q}}}
%\newcommand\C{\ensuremath{\mathbb{C}}}

% Easily typeset systems of equations (French package)
\usepackage{systeme}

% Put x \to \infty below \lim
\let\svlim\lim\def\lim{\svlim\limits}

%Make implies and impliedby shorter
\let\implies\Rightarrow
\let\impliedby\Leftarrow
\let\iff\Leftrightarrow
\let\epsilon\varepsilon
% \let\phi\varphi

% Add \contra symbol to denote contradiction
\usepackage{stmaryrd} % for \lightning
\newcommand\contra{\scalebox{1.5}{$\lightning$}}


% Command for short corrections
% Usage: 1+1=\correct{3}{2}

\definecolor{correct}{HTML}{009900}
\newcommand\correct[2]{\ensuremath{\:}{\color{red}{#1}}\ensuremath{\to }{\color{correct}{#2}}\ensuremath{\:}}
\newcommand\green[1]{{\color{correct}{#1}}}

% horizontal rule
\newcommand\hr{
    \noindent\rule[0.5ex]{\linewidth}{0.5pt}
}

% hide parts
\newcommand\hide[1]{}

% si unitx
\usepackage{siunitx}
\sisetup{}

% Environments
\makeatother
% For box around Definition, Theorem, \ldots
\usepackage{mdframed}
\mdfsetup{skipabove=1em,skipbelow=0em}
\theoremstyle{definition}
\newmdtheoremenv[nobreak=true]{statement}{Statement}[subsection]
\newmdtheoremenv[nobreak=true]{definition}{Definition}[subsection]
%\newmdtheoremenv[nobreak=true]{characteristic}{Characteristic}
%\newmdtheoremenv[nobreak=true]{consequence}{Consequence}
\newmdtheoremenv[nobreak=true]{lemma}{Lemma}[subsection]
\newmdtheoremenv[nobreak=true]{theorem}{Theorem}[subsection]
\newmdtheoremenv[nobreak=true]{proposition}{Proposition}[subsection]
\newmdtheoremenv[nobreak=true]{axiom}{Axiom}[subsection]
\newmdtheoremenv[nobreak=true]{postulate}{Postulate}[subsection]
\newmdtheoremenv[nobreak=true]{prop}{Proposition}[subsection]
\newmdtheoremenv[nobreak=true]{corollary}{Corollary}[subsection]
\newmdtheoremenv{conclusion}{Conclusion}
\newmdtheoremenv{bonus}{Bonus}
\newmdtheoremenv{conjecture}{Conjecture}[subsection]
\newtheorem*{repetition}{Repetition}
\newtheorem*{interlude}{Interlude}
\newtheorem*{notation}{Notation}
\newtheorem*{observation}{Observation}
\newtheorem*{excercise}{Excercise}
\newtheorem*{remark}{Remark}
\newtheorem*{fact}{Fact}
\newtheorem*{claim}{Claim}
\newtheorem*{practice}{Practice}
\newtheorem*{problem}{Problem}
\newtheorem*{terminology}{Terminology}
\newtheorem*{application}{Application}
\newtheorem*{example}{Example}
\newtheorem*{nonexample}{Nonexample}
\newtheorem*{question}{Question}
\newtheorem*{previouslyseen}{As previously seen}

%\newmdtheoremenv[nobreak=true]{definition}{Definition}
%\newtheorem*{notation}{Notation}
%\newtheorem*{remark}{Remark}
%\newtheorem*{note}{Note}
\newtheorem*{observe}{Observe}
\newtheorem*{property}{Property}
\newtheorem*{intuition}{Intuition}
%\newmdtheoremenv[nobreak=true]{theorem}{Theorem}
%\newtheorem{exercise}{Exercise}[section]

% End example and intermezzo environments with a small diamond (just like proof
% environments end with a small square)
\usepackage{etoolbox}
\AtEndEnvironment{example}{\null\hfill$\diamond$}%
\AtEndEnvironment{interlude}{\null\hfill$\diamond$}%

% Fix some spacing
% http://tex.stackexchange.com/questions/22119/how-can-i-change-the-spacing-before-theorems-with-amsthm
\makeatletter
\def\thm@space@setup{%
  \thm@preskip=\parskip \thm@postskip=0pt
}


% Exercise 
% Usage:
% \exercise{5}
% \subexercise{1}
% \subexercise{2}
% \subexercise{3}
% gives
% Exercise 5
%   Exercise 5.1
%   Exercise 5.2
%   Exercise 5.3
\newcommand{\exercise}[1]{%
    \def\@exercise{#1}%
    \subsection*{Exercise #1}
}

\newcommand{\subexercise}[1]{%
    \subsubsection*{Exercise \@exercise.#1}
}


% \lecture starts a new lecture (les in dutch)
%
% Usage:
% \lecture{1}{di 12 feb 2019 16:00}{Inleiding}
%
% This adds a section heading with the number / title of the lecture and a
% margin paragraph with the date.

% I use \dateparts here to hide the year (2019). This way, I can easily parse
% the date of each lecture unambiguously while still having a human-friendly
% short format printed to the pdf.

\usepackage{xifthen}
\def\testdateparts#1{\dateparts#1\relax}
\def\dateparts#1 #2 #3 #4 #5\relax{
    \marginpar{\small\textsf{\mbox{#1 #2 #3 #5}}}
}

\def\@doctype{}%


\newcommand{\lecture}[3]{
    \ifthenelse{\isempty{#3}}{%
        \def\@doctype{Lecture #1}%
    }{%
        \def\@doctype{Lecture #1: #3}%
    }%
    \subsection*{\@doctype}
    \marginpar{\small\textsf{\mbox{#2}}}
}


\newcommand{\lab}[3]{
    \ifthenelse{\isempty{#3}}{%
        \def\@doctype{Lab #1}%
    }{%
        \def\@doctype{Lab #1: #3}%
    }%
    \subsection*{\@doctype}
    \marginpar{\small\textsf{\mbox{#2}}}
}

\newcommand{\homework}[3]{
    \ifthenelse{\isempty{#3}}{%
        \def\@doctype{Homework #1}%
    }{%
        \def\@doctype{Homework #1: #3}%
    }%
    \subsection*{\@doctype}
    \marginpar{\small\textsf{\mbox{#2}}}
}
\newcommand{\assignment}[3]{
    \ifthenelse{\isempty{#3}}{%
        \def\@doctype{Assignment #1}%
    }{%
        \def\@doctype{Assignment #1: #3}%
    }%
    \subsection*{\@doctype}
    \marginpar{\small\textsf{\mbox{#2}}}
}
\newcommand{\studynote}[3]{
    \ifthenelse{\isempty{#3}}{%
        \def\@doctype{Note #1}%
    }{%
        \def\@doctype{Note #1: #3}%
    }%
    \subsection*{\@doctype}
    \marginpar{\small\textsf{\mbox{#2}}}
}

\newcommand{\submit}[3]



% These are the fancy headers
\usepackage{fancyhdr}
\pagestyle{fancy}

% LE: left even
% RO: right odd
% CE, CO: center even, center odd
% My name for when I print my lecture notes to use for an open book exam.
% \fancyhead[LE,RO]{Xiaoyu Liu}

\fancyhead[RO,LE]{\@doctype} % Right odd,  Left even
\fancyhead[RE,LO]{}          % Right even, Left odd

\fancyfoot[RO,LE]{\thepage}  % Right odd,  Left even
\fancyfoot[RE,LO]{}          % Right even, Left odd
\fancyfoot[C]{\leftmark}     % Center

\makeatother




% Todonotes and inline notes in fancy boxes
\usepackage{todonotes}
\usepackage{tcolorbox}

% Make boxes breakable
\tcbuselibrary{breakable}

% Verbetering is correction in Dutch
% Usage: 
% \begin{correction}
%     Lorem ipsum dolor sit amet, consetetur sadipscing elitr, sed diam nonumy eirmod
%     tempor invidunt ut labore et dolore magna aliquyam erat, sed diam voluptua. At
%     vero eos et accusam et justo duo dolores et ea rebum. Stet clita kasd gubergren,
%     no sea takimata sanctus est Lorem ipsum dolor sit amet.
% \end{correction}
\newenvironment{correction}{\begin{tcolorbox}[
    arc=0mm,
    colback=white,
    colframe=green!60!black,
    title=Correction,
    fonttitle=\sffamily,
    breakable
]}{\end{tcolorbox}}

% Noot is note in Dutch. Same as 'verbetering' but color of box is different
\newenvironment{note}[1]{\begin{tcolorbox}[
    arc=0mm,
    colback=white,
    colframe=white!60!black,
    title=#1,
    fonttitle=\sffamily,
    breakable
]}{\end{tcolorbox}}




% Figure support as explained in my blog post.
\usepackage{import}
\usepackage{xifthen}
\usepackage{pdfpages}
\usepackage{transparent}
\newcommand{\incfig}[1]{%
    \def\svgwidth{\columnwidth}
    \import{./figures/}{#1.pdf_tex}
}

% Fix some stuff
% %http://tex.stackexchange.com/questions/76273/multiple-pdfs-with-page-group-included-in-a-single-page-warning
\pdfsuppresswarningpagegroup=1


% My name
\author{Xiaoyu Liu}


% Logseq [[ ]] link support
\newcommand{\logseq}[1]{
	\href{logseq://graph/Feliconut?page=#1}{
	\color{lightgray}{\texttt{[\![\!}}
	\color{cyan}{\textsf{#1}}
	\color{lightgray}{\texttt{\!]\!]}}
}}

\newcommand{\disclaimer}[2]{
    \textbf{Disclaimer.} These are lecture notes taken in #1's course, #2. Correctness of material is not guaranteed. These notes are treated as derived work of #1, and distribution of the material to people unaffliated to the class without the professor's written permission is prohibited.
}
